\frfilename{mti.tex}
\versiondate{21.5.09}
\copyrightdate{2009}





\Loadfont{\twelvebf=cmbx12}

\def\chaptername{Indeks}
\def\sectionname{Topik dan hasil utama}
\def\Krein{Kre\v{\i}n}
\def\med{\mathop{\text{med}}\nolimits}
\def\varinnerprod#1#2{#1\dotproduct#2}
\def\Fn{\mathop{\text{Fn}}\nolimits}

\volumeno=1
\volumeno=2
\volumeno=3
\volumeno=4
\volumeno=5
\ifnum\luluvolumeno>0\volumeno=\luluvolumeno\fi

\ifnum\volumeno<5\def\vfive#1{}\else\def\vfive#1{#1}\fi
\ifnum\volumeno<4\def\vfour#1{}\else\def\vfour#1{#1}\fi
\ifnum\volumeno<3\def\vthree#1{}\else\def\vthree#1{#1}\fi
\ifnum\volumeno<2\def\vtwo#1{}\else\def\vtwo#1{#1}\fi

\def\ind{\mathop{\text{ind}}}

\def\indexiiheader#1{\ifnum\volumeno<2{}\else\indexheader{#1}\fi}
\def\indexiiiheader#1{\ifnum\volumeno<3{}\else\indexheader{#1}\fi}
\def\indexivheader#1{\ifnum\volumeno<4{}\else\indexheader{#1}\fi}
\def\indexvheader#1{\ifnum\volumeno<5{}\else\indexheader{#1}\fi}

\indexmedskip

\centerline{\twelvebf Indeks Jilid 1}





\indexmedskip






\gdef\bottomparagraph{}
\noindent{\bf\sectionname}
   \smallskip
   \let\headlinesectionname=\sectionname

Indeks umum di bawah ini dimaksudkan sebagai indeks yang menyeluruh.   Tidak terelakkan, jumlah entrinya begitu besar sehingga indeks ini sering kali kurang membantu.
Karena itu, saya telah menyiapkan indeks yang lebih ringkas,
dengan anotasi yang lebih baik, yang saya harap dapat membantu pembaca memusatkan perhatian pada bidang tertentu.   Indeks ini tidak mencantumkan definisi, karena entri
bercetak tebal dalam indeks utama diharapkan dapat mengarahkan pembaca secara efisien kepada definisi tersebut;  dan jika Anda tidak menemukan apa yang dicari di sini, tentu saja Anda harus mencoba lagi dalam indeks utama.
Entri yang berbentuk pernyataan matematis sering kali tidak mencantumkan hipotesis penting dan harus dicocokkan dengan pernyataan formal dalam bagian utama karya ini.
































































































\indexmedskip















































































\indexheader{Borel}







himpunan Borel dalam $\BbbR^r$ 111G

----- dan ukuran Lebesgue 114G, 115G, 134F





































\indexmedskip

\indexheader{Cantor}
himpunan dan fungsi Cantor 134G, 134H




\indexheader{\Caratheodory}
konstruksi ukuran dari ukuran luar menurut \Caratheodory~113C

























\indexheader{sentral}




\indexheader{perubahan}
















\indexheader{karakteristik}















































\indexheader{lengkap}


\indexheader{pelengkapan}





\indexheader{konsentrasi}


\indexheader{bersyarat}






\indexheader{konstruksi}
konstruksi ukuran



----- dari ukuran luar (metode \Caratheodory) 113C

----- ukuran subruang 131A






































\indexheader{konvergensi}
teorema-teorema konvergensi (B.Levi, Fatou, Lebesgue) \S123













\indexheader{konveks}




\indexheader{konvolusi}






















\indexheader{terhitung}


himpunan terhitung 111F, 1A1C {\it et seq.}










\indexheader{pencacahan}
ukuran pencacahan 112Bd












\indexheader{Dedekind}















\indexheader{langsung}


























































\indexmedskip


































\indexheader{diperluas}
garis real diperluas \S135













\indexmedskip

\indexheader{Fatou}
Lema Fatou ($\int\liminf\le\liminf\int$
untuk barisan fungsi nonnegatif) 123B




\indexheader{Fej\'er}

























\indexheader{Fourier}


































\indexheader{bebas}








\indexheader{Freese}







\indexheader{Fubini}










\indexheader{fungsi}


\indexheader{fundamental}








































\indexheader{Hahn}






\indexheader{Hardy}


































\indexmedskip








\indexheader{tak tentu}








\indexheader{independen}

















regularitas dalam ukuran



----- (terhadap himpunan kompak) ukuran Lebesgue 134F







\indexheader{integral}










\indexheader{integrasi}
integrasi fungsi bernilai real, konstruksi \S122

----- sebagai fungsional linear positif 122O





----- karakterisasi fungsi terintegralkan 122P, 122R



----- fungsi dan integral dengan nilai dalam $[-\infty,\infty]$ 133A, 135F


















































\indexheader{Kwapien}


\indexmedskip

\indexheader{Lebesgue}












ukuran Lebesgue, konstruksi \S114, \S115

----- sifat-sifat lebih lanjut \S134



Teorema Konvergensi Terdominasi Lebesgue ($\int\lim=\lim\int$
untuk barisan fungsi yang didominasi) 123C

\indexheader{Levi}
Teorema B. Levi ($\int\lim=\lim\int$
untuk barisan fungsi monoton) 123A




\indexheader{pengangkatan}

















\indexheader{Lipschitz}




\indexheader{dapat dilokalkan}




















\indexheader{Loomis}


\indexheader{bawah}

















\indexmedskip

















































\indexheader{terukur}




selubung terukur

----- sifat-sifat elementer 132E



fungsi terukur

----- (bernilai real) \S121

----- ----- jumlah, hasil kali, dan operasi lain atas sejumlah berhingga fungsi 121E

----- ----- limit, infimum, supremum 121F













\indexheader{ukuran}














ruang ukur \S112













\indexheader{monoton}
Teorema Kelas Monoton 136B

\indexheader{monoton}




\indexmedskip



\indexheader{sempit}












\indexheader{tak terukur}
himpunan tak terukur (untuk ukuran Lebesgue) 134B

\indexheader{tak paradoksal}














\indexmedskip

\indexheader{keterurutan}

















\indexheader{Ornstein}


\indexheader{luar}
ukuran luar yang dikonstruksi dari ukuran 132A {\it et seq.}

regularitas luar ukuran Lebesgue 134F










































































































































































































































\indexmedskip





























































































































\indexheader{subruang}
ukuran subruang

----- untuk subruang terukur \S131


































































































































































































































































































\indexmedskip
\indexheader{$\sigma$}
aljabar-$\sigma$ pada himpunan \S111

----- yang dibangkitkan oleh keluarga-keluarga yang diberikan 136B, 136G





\indexheader{$\tau$}



















  

\discrpage

\wheader{Indeks umum}{24}{12}{2}{100pt}


\gdef\newparagraph{}
\gdef\bottomparagraph{}
\def\chaptername{Indeks}
\def\sectionname{Indeks umum}

\noindent{\bf\sectionname}
   \smallskip
   \let\headlinesectionname=\sectionname

\indexmedskip


Rujukan yang dicetak dengan huruf {\bf tebal} mengacu pada definisi;
rujukan dengan huruf {\it miring} merupakan rujukan sepintas.
Definisi yang ditandai dengan $\pmb{>}$ adalah definisi yang pemakaian saya
menyimpang secara berbahaya dari pemakaian sebagian penulis lain.

\indexmedskip



\indexheader{Abel}












































\indexheader{aditif}
fungsional aditif pada suatu aljabar himpunan {\it lihat} aditif berhingga
({\bf 136Xg})




















\indexheader{adjoin}










\indexheader{aljabar}


aljabar himpunan 113Yi, {\bf 136E}, 136F, 136G, 136Xg, 136Xh,
136Xk, 136Ya, 136Yb;  

  {\it lihat juga }
aljabar-$\sigma$ ({\bf 111A})






\indexheader{hampir}







hampir setiap, hampir di mana-mana {\bf 112Dd}








hampir pasti {\bf 112De}

\indexheader{selang-seling}
fungsional selang-seling {\bf 132Yf}
















\indexheader{analitik}
fungsi analitik (kompleks) 133Xc



\indexheader{angelik}


\indexheader{antirantai}




\indexheader{aperiodik}






\indexheader{Archimedean}




































\indexheader{secara asimtotik}











\indexheader{atomik}


\indexheader{tanpa atom}

































\indexheader{aksioma}
aksioma {\it lihat}
  pilihan terhitung



























































\indexmedskip




































\indexheader{bola}


\indexheader{Banach}

































































\indexheader{bidual}




\indexheader{Bienaym\'e}


\indexheader{bilateral}


\indexheader{bilinear}











\indexheader{Bochner}






\indexheader{Boolean}



























\indexheader{Borel}
aljabar Borel {\it lihat} aljabar-$\sigma$ Borel ({\bf 111Gd},
{\bf 135C}) 























fungsi terukur Borel {\bf 121C}, 121D, 121Eg, 121H, 121K, {\bf 121Yd}, 134Fd, 134Xg, 134Yt,
{\bf 135Ef}, 135Xc, 135Xe







himpunan Borel di $\Bbb R$, $\BbbR^r$ {\bf 111G}, {\it 111Yd}, 114G, 114Yd,
115G, 115Yb, {\it 115Yd}, {\it 121Ef}, {\it 121K}, 134F, 134Xd,
{\it 135C}, 136D, 136Xj






aljabar-$\sigma$ Borel (atas subhimpunan $\Bbb R^r$) {\bf 111Gd},
{\it 114Yg-114Yi}, 
121J, 121Xd, 121Xe, 121Yc, 121Yd;

----- (atas ruang lain) {\bf 135C}, 135Xb






\indexheader{berbatas}





himpunan berbatas (di $\BbbR^r$) {\bf 134E}





\indexheader{pembatasan}

















\indexheader{Breiman}






























\indexmedskip















\indexheader{Cantor}
fungsi Cantor {\bf 134H}, 134I



himpunan Cantor {\bf 134G}, 134H, 134I, 134Xf








\indexheader{\Caratheodory}


metode \Caratheodory (untuk mengonstruksi ukuran) 113C, 113D, 113Xa, 113Xd,
113Xg, 113Yc, 113Yk, 114E, 114Xa, 115E, {\it 121Ye},
132Xc, 136Ya
























\indexheader{Carleson}











\indexheader{kategori}


\indexheader{Cauchy}

















\indexheader{selularitas}






















\indexheader{sentral}










\indexheader{rantai}

















\indexheader{karakteristik}
fungsi karakteristik (dari suatu himpunan) {\bf 122Aa}





\indexheader{muatan}


\indexheader{dapat diberi muatan}


\indexheader{Chebyshev}


\indexheader{pilihan}
pilihan, aksioma 134C, 1A1G; 
 {\it lihat juga} pilihan terhitung
















\indexheader{lingkaran}




\indexheader{buka-tutup}



\indexheader{tertutup}






interval tertutup (di $\Bbb R$ atau $\BbbR^r$) 114G, 115G, 1A1A

himpunan tertutup (di ruang topologis) 134Fb, 134Xd, {\bf 1A2E}, 1A2F,
1A2G





\indexheader{penutupan}


\indexheader{club (tertutup tak terbatas)}


\indexheader{gugus}






\indexheader{koarea}





























































\indexheader{kofinit}




















\indexheader{komutatif}






\indexheader{komutator}




\indexheader{saling komutatif}


\indexheader{kompak}


























\indexheader{kompaktifikasi}














\indexheader{komplemen}


\indexheader{komplementer}


\indexheader{terkomplemen}





\indexheader{lengkap}











ukuran (ruang ukur) lengkap {\bf 112Df}, 113Xa, 122Ya
















\indexheader{dilengkapi}


\indexheader{sepenuhnya}
















\indexheader{pelengkapan}


























\indexheader{kompleksifikasi}


\indexheader{bernilai kompleks}
fungsi bernilai kompleks \S133

\indexheader{komponen}
komponen 
  (dalam ruang topologis) {\it 111Ye}







































\indexheader{koterabaikan}
himpunan koterabaikan {\bf 112Dc}

\indexheader{terhubung}







\indexheader{konjugat}










\indexheader{konsisten}

















\indexheader{kontinu}
fungsi kontinu 121D, 121Yd

















\indexheader{kontinuum}


























\indexheader{konvergen}






{barisan konvergen }
{\bf 135D}



\indexheader{konveks}




































\indexheader{terhitung}
himpunan terhitung {\bf 111F}, 114G, 115G, \S1A1







pilihan terhitung, aksioma {\it 134C}























































\indexheader{pencacahan}
ukuran cacah {\bf 112Bd}, {\it 122Xd}, 122 {\it catatan}








\indexheader{penutup}
penutup {\it lihat} selubung terukur ({\bf 132D})






















\indexheader{kozero}




\indexheader{siklus}


\indexheader{siklik}


\indexheader{silinder}






\indexmedskip



\indexheader{Davies}




\indexheader{De Morgan}




\indexheader{desimal}







\indexheader{dekomposisi}




\indexheader{menurun}




\indexheader{Dedekind}


{himpunan terurut lengkap Dedekind }
{\it 135Ba}





















































\indexheader{turunan}
turunan suatu fungsi\ {\it lihat} turunan parsial












\indexheader{determinan}


\indexheader{ditentukan}







\indexheader{Setan}
Tangga Setan {\it lihat} fungsi Cantor ({\bf 134H})











\indexheader{keterdiferensialan}


\indexheader{dapat didiferensialkan}
fungsi yang dapat didiferensialkan (satu variabel) 123D



\indexheader{mendiferensialkan}
diferensiasi di bawah tanda integral 123D

\indexheader{tersebar}


\indexheader{dilatasi}










\indexheader{Dirac}
ukuran Dirac {\bf 112Bd}

\indexheader{langsung}
citra langsung (dari suatu himpunan oleh fungsi atau relasi) 1A1B


























\indexheader{disintegrasi}
disintegrasi suatu ukuran 123Ye


\indexheader{saling lepas}
keluarga himpunan yang saling lepas {\bf 112Bb}









\indexheader{distribusi}










\indexheader{distributif}








\indexheader{divergensi}












\indexheader{terdominasi}
Teorema Konvergensi Terdominasi {\it lihat} Teorema Konvergensi Terdominasi Lebesgue
({\bf 123C})




\indexheader{Doob}


\indexheader{ganda}


\indexheader{berganda}






































\indexheader{dual}



















\indexheader{Dunford}




\indexheader{diadik}








\indexheader{Dye}





\indexheader{Dynkin}
kelas Dynkin {\bf 136A}, 136B, 136Xb




\indexmedskip

\indexheader{Eberlein}












\indexheader{Effros}




\indexheader{Egorov}
teorema Egorov 131Ya





\indexheader{vektor eigen}









\indexheader{dapat dibenamkan}


\indexheader{pembenaman}


\indexheader{empiris}















\indexheader{entropi}




\indexheader{mengenumerasi}


\indexheader{selubung}
selubung {\it lihat}
selubung terukur ({\bf 132D})




















  

















\indexheader{ekuivalen}
ekuivalen {\bf 121B}



































\indexheader{secara esensial}











\indexheader{Euklides}


topologi Euklides \S1A2



\indexheader{genap}




\indexheader{dapat dipertukarkan}


\indexheader{pertukaran}


\indexheader{eksauksi}

















\indexheader{secara eksponensial}


\indexheader{eksponensiasi}


\indexheader{diperluas}


garis real diperluas {\it 121C}, \S135

\indexheader{perluasan}


perluasan fungsional aditif berhingga 113Yi



perluasan ukuran 113Yc, 113Yh, 132Yd

















\indexmedskip















\indexheader{Fatou}
Lema Fatou {\bf 123B}, 133Kb, 135Gb, 135Hb









\indexheader{Federer}


\indexheader{Fej\'er}









\indexheader{Feller}













\indexheader{medan}
medan (himpunan) {\it lihat} aljabar ({\bf 136E})
















\indexheader{halus}







\indexheader{berhingga}


















\indexheader{secara berhingga}
fungsi (atau fungsional) aditif berhingga pada suatu aljabar himpunan 136Xg, 136Ya,
136Yb



\vindexheader{secara berhingga}{36}























































\indexheader{Fourier}




deret Fourier {\it 121G}



transformasi Fourier {\bf 133Xd}, {\bf 133Yc}













\indexheader{Fr\'echet}






\indexheader{bebas}



















\indexheader{Fremlin}




\indexheader{Frol\'\i k}


\indexheader{Frostman}











\indexheader{penuh}




ukuran luar penuh {\bf 132F}, 132Yd, 134D, 134Yt











\indexheader{fungsi}
fungsi 1A1B



\indexheader{fundamental}





\indexmedskip




\indexheader{Gaifman}








\indexheader{permainan}


\indexheader{gamma}











\indexheader{tolok}




\indexheader{Gauss}




\indexheader{Gaussian}











\indexheader{dibangkitkan}






aljabar himpunan yang dibangkitkan (atau aljabar-$\sigma$\nobreak yang dibangkitkan) {\bf 111Gb}, 111Xe, 111Xf,
{\it 121J}, 121Xd, 121Yc, 136B, 136C, 136G, 136H, 136Xc, 136Xk,
136Yb, 136Yc



\indexheader{pembangkit}







\indexheader{Geschke}
























\indexheader{G\"odel}







\indexheader{gradien}


\indexheader{graf}





\indexheader{Green}




\indexheader{grup}












\indexmedskip









































\indexheader{setengah}
interval setengah terbuka (di $\Bbb R$ atau $\BbbR^r$) {\bf 114Aa}, 114G,
{\it 114Yj}, {\bf 115Ab}, 115Xa, {\it 115Xc}, {\it 115Yd}


















































































\indexheader{secara herediter}

















\indexheader{Hilbert}




































\indexheader{homomorfik}


\indexheader{homomorfisme}








\indexheader{selubung}


\indexheader{Humpty}

















\indexmedskip

\indexheader{ideal}
ideal dalam suatu aljabar himpunan

  {\it lihat} ideal-$\sigma$ ({\bf 112Db})











\indexheader{idempoten}







\indexheader{citra}


ukuran citra {\bf 112Xf}, 123Ya, 132Yb





























\indexheader{tak tentu}


integral tak tentu 131Xa





























































\indexheader{tak hingga}
tak hingga 112B, 133A, \S135

\indexheader{informasi}









\indexheader{injektif}




\indexheader{dalam}


ukuran dalam 113Yg

----- ----- yang didefinisikan oleh suatu ukuran 113Yh 






\indexheader{regular dalam}















\indexheader{terintegralkan}
fungsi terintegralkan \S122 ({\bf 122M}), {\it 123Ya}, {\bf 133B}, {\bf
133Db}, 133F, 133Jd, 133Xa, {\it 135Fa}

\indexheader{integral}
integral \S122 ({\bf 122E}, {\bf 122K}, {\bf 122M}), {\bf 133B}, {\bf
133D}; 
  {\it lihat juga}
fungsi terintegralkan, integral Lebesgue ({\bf 122Nb}),
integral bawah ({\bf 133I}), integral Riemann ({\bf 134K}), integral atas ({\bf 133I})














\indexheader{interior}










\indexheader{irisan}


\indexheader{interval}
interval
  {\it lihat}
interval setengah terbuka ({\bf 114Aa}, {\bf 115Ab}),
interval terbuka ({\bf 111Xb})



\indexheader{invarian}


















 

\indexheader{balik}




citra balik (dari suatu himpunan oleh fungsi atau relasi) {\bf 1A1B}


fungsi pelestari ukuran melalui citra balik
134Yl-134Yn



\indexheader{inversi}




\indexheader{dapat diinvers}


\indexheader{involusi}












\indexheader{isodiametrik}


\indexheader{terisolasi}

























\indexheader{Jacobian}













































\indexheader{Kawada}












































\indexheader{Koml\'os}











\indexheader{\Krein}




\indexheader{Kronecker}


\indexheader{Kuratowski}









\indexheader{Kwapien}




\indexmedskip

\indexheader{Lacey}


\indexheader{Laplace}


transformasi Laplace {\bf 123Xc}, {\it 123Yb}, {\bf 133Xd}




\indexheader{secara lateral}


\indexheader{kekisi}








\indexheader{hukum}








\indexheader{Lebesgue}
Lebesgue, H.\   Jilid\ 1 {\it pendahuluan}











Teorema Konvergensi Terdominasi Lebesgue 123C, 133G



fungsi terintegralkan secara Lebesgue {\bf 122Nb}, 122Yb, 122Ye, {\it
122Yf}

integral Lebesgue {\bf 122Nb}

fungsi terukur Lebesgue {\bf 121C}, 121D, 134Xg, 134Xj

himpunan terukur Lebesgue {\bf 114E}, 114F, 114G, {\it 114Xe}, 114Ye, {\bf
115E}, 115F, 115G

ukuran Lebesgue (pada $\Bbb R$) \S114 ({\bf 114E}), 131Xb, 133Xd, 133Xe,
134G-134L


----- ----- (pada $\BbbR^r$) \S115 ({\bf 115E}), 132C, {\it 132Ef}, 133Yc,
\S134




----- ----- (pada subhimpunan lain dari $\BbbR^r$) {\bf 131B}





himpunan terabaikan Lebesgue {\bf 114E}, {\bf 115E}, 134Yk



ukuran luar Lebesgue {\bf 114C}, 114D, 114Xc, {\it 114Yd}, {\bf
115C}, 115D, 115E, 115Xb, 115Xd, 115Yb, 132C, 134A, 134D, 134Fa



ukuran Lebesgue-Stieltjes {\bf 114Xa}, {\it 114Xb}, 114Yb, 114Yc, 114Yf, 131Xc,
132Xg, 134Xd



\indexheader{kiri}








\indexheader{panjang}




panjang suatu interval {\bf 114Ab}




\indexheader{Levi}
teorema B. Levi {\bf 123A}, 123Xa, 133Ka, 135Ga, 135Hb



\indexheader{L\'evy}








\indexheader{leksikografis}




\indexheader{Liapounoff}




\indexheader{pengangkatan}








\indexheader{limit}















\indexheader{Lindel\"of}






\indexheader{linear}




\ifdim\pagewidth>467pt\fi

\fontdimen3\tenrm=1.67pt












\indexheader{tertaut}






\indexheader{Lipschitz}






\indexheader{lokal}




















































\indexheader{Lo{\grv e}ve}


\indexheader{logistik}


\indexheader{Loomis}


\indexheader{Lorentz}




\indexheader{bawah}






integral bawah {\bf 133I}, 133J, 133Xf, 133Yd, 135Ha





integral Riemann bawah {\bf 134Ka}








\indexheader{Lusin}




teorema Lusin 134Yd

\indexheader{Luxemburg}




\indexmedskip











































































\indexheader{martingal}











\indexheader{aliran maksimum}


\indexheader{maksimal}










\indexheader{Mazur}




\indexheader{McMillan}


\indexheader{McShane}




\indexheader{kurus}




\indexheader{rataan}






\indexheader{terukur}






penutup terukur {\it lihat} selubung terukur ({\bf 132D})

selubung terukur {\bf 132D}, 132E, 132F, 132Xg, 134Fc, 134Xd;  
  {\it lihat juga }
  berukuran luar penuh ({\bf 132F})



fungsi terukur (bernilai dalam $\Bbb R$) \S121 ({\bf 121C}),
122Ya

----- -----  (bernilai dalam $\BbbR^r$) {\bf 121Yd}

----- -----  (bernilai dalam ruang lain) {\bf 133Da}, 133E, 133Yb,
{\bf 135E}, {\it 135Xd}, 135Yf  

----- -----  ($(\Sigma,\Tau)$-terukur) {\bf 121Yc}

----- -----  {\it lihat juga} terukur Borel, terukur Lebesgue






himpunan terukur {\bf 112A};
  {\it lihat juga }
terukur secara relatif ({\bf 121A})


ruang terukur {\bf 111Bc}



transformasi terukur
  {\it lihat} fungsi pelestari ukuran melalui citra balik

\indexheader{ukuran}
ukuran {\bf 112A}

----- (dalam `$\mu$ mengukur $E$', `$E$ diukur oleh $\mu$') {\bf $\pmb{>}$112Be}














































\indexheader{ruang ukur}
ruang ukur \S112 ({\bf 112A})


















\indexheader{metrik}








\indexheader{secara metrik}







































\indexheader{termoderasi}




\indexheader{termodifikasi}


\indexheader{modular}




\indexheader{modulasi}


\indexheader{monokompak}


\indexheader{monoton}
Teorema Kelas Monoton 136B-136D, 
136Xc, 136Xf

Teorema Konvergensi Monoton {\it lihat} Teorema B. Levi (123A)



\indexheader{monoton}
fungsi monoton 121D








\indexheader{multilinear}














\indexmedskip

\indexheader{Nachbin}


\indexheader{Nadkarni}





















\indexheader{sempit}


\indexheader{alami}









\indexheader{terabaikan}
himpunan terabaikan {\bf 112D}, 131Ca; 
  {\it lihat juga} 
terabaikan Lebesgue ({\bf 114E}, {\bf 115E}), ideal nol ({\bf 112Db})




\indexheader{lingkungan}







\indexheader{jaringan}




\indexheader{Neumann}













\indexheader{Nikod\'ym}






\indexheader{tidak}


\indexheader{tak-menurun}
barisan tak-menurun dari himpunan {\bf 112Ce}

\indexheader{tak-menaik}
barisan tak-menaik dari himpunan {\bf 112Cf}

\indexheader{tak terukur}
himpunan tak terukur 134B, 134D, 134Xc

\indexheader{tak-utama}













\indexheader{norma}


\indexheader{normal}




























\indexheader{ternormalisasi}






\indexheader{penormal}


\indexheader{bernorma}














\indexheader{tak di mana pun}














\indexheader{nol}
ideal nol suatu ukuran {\bf 112Db}




himpunan nol {\it lihat} terabaikan ({\bf 112Da})


\indexmedskip

\indexheader{ganjil}







\indexheader{Ogasawara}
















\indexheader{terbuka}
selang terbuka {\it 111Xb}, 114G, 115G, 1A1A



himpunan terbuka (dalam $\BbbR^r$) {\bf 111Gc}, {\it 111Yc}, {\it 114Yd}, {\it 115G},
{\it 115Yb}, {\bf 133Xc}, 134Fa, 134Xe,
{\bf 135Xa}, {\bf 1A2A}, 1A2B, 1A2D;
  (dalam $\Bbb R$) 111Gc, 111Ye, {\it 114G}, 134Xd



















































































































































\indexheader{luar}




\indexheader{ukuran luar}
ukuran luar \S113 ({\bf 113A}), 114Xd, {\bf 132B}, {\it 132Xg},
136Ya;  
  {\it lihat juga} ukuran luar Lebesgue ({\bf 114C}, {\bf 115C}),
  ukuran luar reguler ({\bf 132Xa})

----- ----- yang didefinisikan dari suatu ukuran 113Ya,
132A-132E ({\bf 132B}), 
132Xa-132Xi, 
132Xk, 132Ya-132Yc, 
133Je














\indexmedskip
















\indexheader{parsial}
turunan parsial {\it 123D}







\indexheader{secara parsial}




\indexheader{partisi}
partisi {\bf 1A1J}





\indexheader{Peano}
kurva Peano 134Yl-134Yo 

\indexheader{sempurna}








\indexheader{secara sempurna}


\indexheader{keliling}




\indexheader{periodik}












\indexheader{Perron}


\indexheader{Pettis}






\indexheader{Pfeffer}


















\indexheader{Plancherel}







\indexheader{Poincar\'e}


\indexheader{terhitung-titik}


\indexheader{berhingga-titik}




\indexheader{bertumpu pada titik}
ukuran yang bertumpu pada titik {\bf 112Bd};\ 
  {\it lihat juga} ukuran Dirac ({\bf 112Bd})








































































{himpunan kuasa $\Cal P\Bbb N$ }
1A1Hb



















\indexheader{dapat diprediksi}





\indexheader{pra-Radon}


\indexheader{presque}
presque partout {\bf 112De}




\indexheader{primitif}




























\indexheader{produk}





  



























































































\indexheader{pseudo-sederhana}
fungsi pseudo-sederhana {\bf 122Ye}, 133Ye



\indexheader{tarikan-balik}


\indexheader{secara murni}













\indexmedskip

\indexheader{kuasi-diadik}


\indexheader{kuasi-homogen}





\indexheader{kuasi-rapat-orde}







\indexheader{kuasi-sederhana}
fungsi kuasi-sederhana {\bf 122Yd}, 133Yd



























\indexmedskip


























































































\indexheader{rekuren}



































































ukuran luar reguler 132C, {\bf 132Xa}





\indexheader{secara reguler}






\indexheader{relasi}
relasi 1A1B












\indexheader{secara relatif}




















himpunan terukur secara relatif {\bf 121A}












\indexheader{berulang}


\indexheader{representasi}

















\indexheader{mempertahankan}




\indexheader{retrak}




\indexheader{balik}







\indexheader{Riemann}


fungsi terintegralkan Riemann {\bf 134K}, 134L

integral Riemann {\bf 134K}





\indexheader{Riesz}































































\indexmedskip

\indexheader{Saks}






\indexheader{saltus}













\indexheader{secara skalar}


\indexheader{terpencar}





\indexheader{Schr\"oder}


\indexheader{Schwartz}


\indexheader{ala Schwartz}




\indexheader{kedua}














\indexheader{adjoin-diri}


\indexheader{menopang-diri}







\indexheader{semikompak}


\indexheader{semikontinu}


\indexheader{semi-definit}


\indexheader{semi-berhingga}






\indexheader{semigrup}


\indexheader{semi-martingal}


\indexheader{seminorma}


\indexheader{semi-radon}


\indexheader{semiring}
semiring himpunan {\bf 115Ye}



\indexheader{separabel}











\indexheader{memisahkan}


\indexheader{terpisah}





\indexheader{pemisahan}














\indexheader{secara sekuensial}
































































\indexheader{Sierpi\'nski}
Teorema Kelas Sierpi\'nski {\it lihat} Teorema Kelas Monoton (136B)



\indexheader{bertanda}











\indexheader{sederhana}
fungsi sederhana \S122 ({\bf $\pmb{>}$122A})











\indexheader{Sina\v\i}










\indexheader{miring}




\indexheader{slalom}




\indexheader{kecil}

















\indexheader{solid}































\indexheader{ruang}
kurva pengisi ruang 134Yl




\indexheader{spektral}


\indexheader{spektrum}






















\indexheader{persegi}






\indexheader{penstabil}






  

























\indexheader{Stieltjes}
ukuran Stieltjes {\it lihat}
ukuran Lebesgue-Stieltjes ({\bf 114Xa})










\indexheader{secara stokastik}


\indexheader{Stone}














\indexheader{penghentian}


\indexheader{langsung}


\indexheader{Strassen}


\indexheader{strategi}












\indexheader{kuat}














































\indexheader{subdivisi}


\indexheader{subgrup}





\indexheader{subhomomorfisme}


\indexheader{subkisi}


\indexheader{submartingal}


\indexheader{subukuran}


\indexheader{submodular}


\indexheader{subgelanggang}




\indexheader{subruang}
ukuran subruang 113Yb  

----- -----  pada subhimpunan terukur 131A, {\bf 131B}, 131C, 132Xb,
135I  

----- -----  (integrasi terhadap ukuran subruang) {\bf 131D},
131E-131H, 
131Xa-131Xc, 
{\it 133Dc}, 133Xb, 135I







aljabar-$\sigma$ subruang {\bf 121A}












\indexheader{jumlah}
jumlah atas himpunan indeks sembarang 112Bd

jumlah ukuran 112Yf, 122Xi




\indexheader{dapat dijumlahkan}








\indexheader{supermodular}


\indexheader{tumpuan}










\indexheader{tertumpu}


{tertumpu {\it lihat}}\
tertumpu-titik ({\bf 112Bd})
































\indexmedskip























\indexheader{bertemper}






\indexheader{tensor}


\indexheader{tenda}


\indexheader{uji}




\indexheader{tebal}
himpunan tebal {\it lihat} berukuran luar penuh ({\bf 132F})









\indexheader{Tietze}


\indexheader{ketat}









































\indexheader{topologi}














\indexheader{secara total}











himpunan terurut total {\it 135Ba}






\indexheader{jejak}




{jejak }
(dari suatu aljabar-$\sigma$) {\it lihat} aljabar-$\sigma$ subruang ({\bf 121A})









\indexheader{transitif}


\indexheader{invarian terhadap translasi}










{invarian terhadap translasi }
ukuran 114Xf, 115Xd, 134A, 134Ye, 134Yf





\indexheader{transportasi}


\indexheader{transversal}








\indexheader{secara sejati}



\indexheader{terpangkas}






















\indexheader{Tychonoff}


\indexheader{tipe}



\indexmedskip









\indexheader{ultrafilter}








\indexheader{belum dilengkapi}






\indexheader{uniferent}


\indexheader{seragam}












































































































\indexheader{penyilangan-naik}








\indexheader{atas}








integral atas {\bf 133I}, 133J-133L, 
133Xf, 133Yd, {\bf 135H}, 135I, 135Yc

integral Riemann atas {\bf 134Ka}






























\indexheader{Urysohn}






\indexheader{biasa}









\indexmedskip










\indexheader{variasi}






\indexheader{vektor}








\indexheader{sangat}




\indexheader{Vietoris}


\indexheader{secara virtual}
fungsi terukur secara virtual {\bf 122Q}, 122Xe, 122Xf, 135Ia

\indexheader{Vitali}
konstruksi Vitali atas suatu himpunan tak terukur 134B









\indexheader{volume}
volume {\bf 115Ac}






\indexheader{Wald}


\indexheader{mengembara}




\indexheader{lemah}





























































\indexheader{Weierstrass}




































\indexheader{Wirtinger}














\indexheader{Young}





























\indexmedskip 

\indexheader{$\scriptstyle\pmb{A}$}















a.e.\ (`hampir di mana-mana') {\bf 112Dd}

a.s.\ (`hampir pasti') {\bf 112De}



























\indexmedskip
\indexheader{$\scriptstyle\pmb{C}$}











$\Bbb C$ = himpunan bilangan kompleks























































\indexmedskip
\indexheader{$\scriptstyle\pmb{D}$}

























$\diam$ (dalam $\diam A$) = diameter



$\dom$ (dalam $\dom f$):  domain suatu fungsi $f$


\indexheader{$\scriptstyle\pmb{E}$}








\indexheader{$\scriptstyle\pmb{F}$}






















\indexheader{$\scriptstyle\pmb{G}$}










\indexheader{$\scriptstyle\pmb{H}$}














































\indexmedskip
\indexheader{$\scriptstyle\pmb{L}$}

























$\eusm L^0$ (dalam $\eusm L^0(\mu)$) 121Xb
















$\eusm L^1$ (dalam $\eusm L^1(\mu)$) {\it 122Xc}














































$\liminf$ (dalam $\liminf_{n\to\infty}$) \S1A3 ({\bf 1A3Aa});
  (dalam $\liminf_{t\downarrow 0}$) {\bf 1A3D}

$\limsup$ (dalam $\limsup_{n\to\infty}$) \S1A3 ({\bf 1A3Aa});
  (dalam $\limsup_{t\downarrow 0}$) {\bf 1A3D}
















































































\indexmedskip
\indexheader{$\scriptstyle\pmb{N}$}



$\Bbb N\times\Bbb N$ 111Fb























\indexmedskip
\indexheader{$\scriptstyle\pmb{P}$}



$\Cal P$ {\it lihat} himpunan kuasa

























p.p.\ (`presque partout') {\bf 112De}



\indexmedskip
\indexheader{$\scriptstyle\pmb{Q}$}

$\Bbb Q$ (himpunan bilangan rasional) 111Eb, 1A1Ef





\indexmedskip
\indexheader{$\scriptstyle\pmb{R}$}

$\Bbb R$ (himpunan bilangan real) 111Fe, 1A1Ha






$\overline{\Bbb R}$ {\it lihat} garis real diperluas (\S135)






















\indexheader{$\scriptstyle\pmb{S}$}






































\indexheader{$\scriptstyle\pmb{T}$}































\indexmedskip
\indexheader{$\scriptstyle\pmb{U}$}

$U$ (dalam $U(x,\delta)$) {\bf 1A2A}













\indexheader{$\scriptstyle\pmb{V}$}






\indexheader{$\scriptstyle\pmb{W}$}














\indexmedskip
\indexheader{$\scriptstyle\pmb{Z}$}

$\Bbb Z$ (himpunan bilangan bulat) 111Eb, 1A1Ee







ZFC {\it lihat} teori himpunan Zermelo-Fraenkel

 



 






 













































\indexmedskip
\indexheader{$\scriptstyle\pmb{\pi}$}







Teorema $\pi$-$\lambda$ {\it lihat} Teorema Kelas Monoton (136B)

\indexmedskip





\indexmedskip
\indexheader{$\scriptstyle\pmb{\sigma}$}







aljabar-$\sigma$ himpunan {\bf 111A}, 111B,
111D-111G, 
111Xc-111Xf, 
111Yb, 136Xb, 136Xi; 
  {\it lihat juga}
aljabar-$\sigma$ Borel ({\bf 111G})
















medan-$\sigma$ {\it lihat} aljabar-$\sigma$ ({\bf 111A})













\indexheader{ideal-$\sigma$}


{ideal-$\sigma$ }
(himpunan) {\bf 112Db}


\indexheader{$\sigma$-terisolasi}


\indexheader{$\sigma$-tertaut}




\indexheader{$\sigma$-diskret-metrik}


\indexheader{$\sigma$-lengkap-orde}


\indexheader{$\sigma$-kontinu-orde}





\indexheader{subaljabar-$\sigma$}




\indexheader{$\sigma$-subhomomorfisme}


\indexheader{$(\sigma,\infty)$-distributif}


\indexmedskip
\indexheader{$\scriptstyle\pmb{\Sigma}$}







$\sum_{i\in I}a_i$ {\bf 112Bd}





































\indexmedskip
\indexheader{$\scriptstyle\pmb{\chi}$}

$\chi$ (dalam $\chi A$, dengan $A$ suatu himpunan) {\bf 122Aa}






\indexheader{$\scriptstyle\pmb{\omega}$}


































\indexmedskip
\indexheader{simbol khusus}

$\setminus$ (dalam $E\setminus F$, `selisih himpunan') {\bf 111C}

$\symmdiff$ (dalam $E\symmdiff F$, `selisih simetris') {\bf 111C}







$\bigcup$ (dalam $\bigcup_{n\in\Bbb N}E_n$) {\bf 111C};
  (dalam $\bigcup\Cal A$) {\bf 1A1F}

$\bigcap$ (dalam $\bigcap_{n\in\Bbb N}E_n$) {\bf 111C};
  (dalam $\bigcap\Cal E$) {\bf 1A2F}

$\vee$, $\wedge$ (dalam suatu kisi) 121Xb





$\restr$ (dalam $f\restr A$, pembatasan suatu fungsi pada suatu himpunan) {\bf
121Eh}





$\eae$ {\bf 112Dg}, 112Xe

$\leae$ {\bf 112Dg}, 112Xe

$\geae$ {\bf 112Dg}, 112Xe






























*
 
 
(dalam $\mu^*$) {\it lihat} ukuran luar yang diturunkan dari suatu ukuran ({\bf 132B})

$_*$ (dalam $\mu_*$) {\it lihat} ukuran dalam yang diturunkan dari suatu ukuran
(113Yh)
















$\int$ (dalam $\int f$, $\int fd\mu$, $\int f(x)\mu(dx)$) {\bf 122E}, {\bf
122K}, {\bf 122M}, 122Nb; 
  {\it lihat juga} integral atas, integral bawah ({\bf 133I})



----- (dalam $\int_Af$) {\bf 131D};  
  {\it lihat juga} ukuran subruang



$\overline{\int}$ {\it lihat} integral atas ({\bf 133I})

$\underline{\int}$ {\it lihat} integral bawah ({\bf 133I})





$\Rint$ {\it lihat} integral Riemann ({\bf 134K})
































































$^+$
  
  
  (dalam $f^+$, dengan $f$ suatu fungsi) {\bf 121Xa}
  
  
  


$^-$ (dalam $f^-$, dengan $f$ suatu fungsi) {\bf 121Xa}






















































$\infty$ {\it lihat} tak hingga





$[\hskip1em]$ (dalam $[a,b]$) {\it lihat} selang tertutup ({\bf 115G},
{\bf 1A1A});
  (dalam $f[A]$, $f^{-1}[B]$, $R[A]$, $R^{-1}[B]$) {\bf 1A1B}





$\coint{\hskip1em}$ (dalam $\coint{a,b}$) {\it lihat} selang setengah-terbuka
({\bf 115Ab}, {\bf 1A1A})

$\ocint{\hskip1em}$ (dalam $\ocint{a,b}$) {\it lihat} selang setengah-terbuka
({\bf 1A1A})

$\ooint{\hskip1em}$ (dalam $\ooint{a,b}$) {\it lihat} selang terbuka
({\bf 115G}, {\bf 1A1A})

















\frnewpage

